% -*- coding:utf-8 -*-
\documentclass[12 pt, a4paper, twoside]{article}

\usepackage[utf8]{inputenc}
\usepackage[T1]{fontenc}
\usepackage[spanish, activeacute]{babel}
\usepackage[pdftex, bookmarks, hyperfootnotes=false, colorlinks=true,%
            urlcolor=blue, linkcolor=black,  citecolor=black,%
            pagecolor=black, anchorcolor=black, breaklinks=true]{hyperref}
\usepackage{times}
\usepackage{eurosym}

\usepackage[pdftex]{graphicx}
%\usepackage{rotating}
\graphicspath{{img/}}

\usepackage{epstopdf}
\epstopdfsetup{outdir=img/,
  suffix=-generated}

\usepackage[pdftex]{color}
\definecolor{gray97}{gray}{.97}
\definecolor{gray75}{gray}{.75}
\definecolor{gray45}{gray}{.45}

\usepackage{listingsutf8}
\lstset{
  inputencoding = utf8/latin1,
                  %
                  frame = Ltb,
                  framerule = 0 pt,
                  aboveskip = 0.5 cm,
                  framextopmargin = 3 pt,
                  framexbottommargin = 3 pt,
                  framexleftmargin = 0.4 cm,
                  framesep = 0 pt,
                  rulesep = .4 pt,
                  backgroundcolor = \color{gray97},
                  rulesepcolor = \color{black},
                  %
                  stringstyle = \ttfamily,
                  showstringspaces  =  false,
                  basicstyle = \small\ttfamily,
                  commentstyle = \color{gray45},
                  keywordstyle = \bfseries,
                  %
                  numbers = left,
                  numbersep = 15 pt,
                  numberstyle = \tiny,
                  numberfirstline  =  false,
                  breaklines = true
}

\hypersetup{colorlinks = true, urlcolor = blue,
            pdftitle = Game Design Document,
            pdfsubject = Práctica de Ingeniería del Software II,
            pdfauthor = Varios}

\hoffset -0.54 cm
\voffset -0.54 cm
\textheight = 22 cm
\textwidth = 17 cm
\topmargin = 0 cm
\oddsidemargin = 0 cm
\evensidemargin = 0 cm
%\parindent = 0 cm
\parskip = 6 pt

\usepackage{fancyhdr}
\pagestyle{fancy}
\fancyhf[EH,EF,OH,OF]{}
\fancyhf[OLH]{Práctica 2}
\fancyhf[ERH]{Ingeniería del Software 2011-12}
\fancyhf[ELH]{\thepage}
\fancyhf[ORH]{\thepage}

\pdfimageresolution = 300

\title{Práctica 2: {\em Game Design Document}\\Ingeniería del Software 2011-12}
\author{Manuel José Abaldea García-Pliego\\
Luis Miguel Garcia-Muñoz Pérez\\
Eduardo Monroy Martínez\\
Felipe Terriza García-Muñoz}
\date{}

\begin{document}

\maketitle

%\begin{titlepage}
\begin{center}

\begin{figure}[h]
\centering
\includegraphics[width = 12 cm]{esi_color.eps}
\end{figure}

{\huge Ingeniería del Software II 2011-12.\\}

{\huge Desarrollo de una aplicación en Android.\\}
\end{center}

\begin{center}
29 de Abril del 2012
\end{center}

\begin{center}
{\large Manuel José Abaldea García-Pliego

Luis Miguel Garcia-Muñoz Pérez

Eduardo Monroy Martínez

Felipe Terriza García-Muñoz}
\end{center}

\end{titlepage}


%\newpage

%\tableofcontents

\section{Cambios}

\begin{itemize}
  \item 2 de Marzo del 2012: Primera versión del documento.
\end{itemize}

\section{Introducción}

En este documento describimos el videojuego que vamos a crear. Es de carácter
iterativo, de manera que se obtendrá una nueva versión de él al final de cada
ciclo de desarrollo. La versión final será aquella que vaya junto con la
versión definitiva del videojuego.

%\input{genero}

\section{Jugadores}

El juego es para un jugador por dispositivo y sin juego en red.

%\input{historia}
%\input{bocetos} (se añadirá en futuras iteraciones)
%\input{lookandfeel}
%\input{interfaz}
%\input{objetivos}
%\input{reglas}
%\input{características}
%\input{gameplay}
%\input{niveles}
%\input{requerimientos}

\section{Marketing}

% -*- coding:utf-8 -*-
\section{Presupuesto}

El coste estimado para los 4 meses de desarrollo es el siguiente:

\begin{tabular}
  Concepto & Precio unitario & Total \\
  Salarios & 18000 \euro/persona & 18000 \euro \\
  Internet & 120 \euro/persona & 480 \euro \\
  Desplazamientos & 120 \euro/persona & 480
\end{tabular}

{\bf Coste total:} 19020 \euro

Fondo para costes imprevistos: 5000 \euro


\section{Trabajo pendiente}

Con la finalización de la primera versión de éste documento, se da por acabada
la fase de pre-producción, iniciándose la fase de producción y comenzando con
el desarrollo propiamente dicho del videojuego.

\end{document}
